Un câble coaxial est formé d'un conducteur central plein de rayon $a = \qty{\v{a}}{mm}$, parcouru par un courant de \qty{\v{I}}{A} uniformément réparti dans sa section, et d'une gaine conductrice mince de rayon $b = \qty{\v{b}}{mm}$ qui ramène le même courant en sens inverse. ($\mu_0 = 4\pi \times \qty{e-7}{T.m/A}$)

\begin{enumerate}[label=\alph*)]
    \item À l'aide du théorème d'Ampère, déterminez le champ magnétique à \qty{\v{r1}}{mm} de l'axe (dans le conducteur central).
    \item Déterminez le champ magnétique à \qty{\v{r2}}{mm} de l'axe (entre le conducteur central et la gaine).
    \item Que vaut le champ magnétique à l'extérieur du câble ($r > b$) ? Justifiez.
\end{enumerate}
